\solchapter{geo}

\sol{geo}{1} $g = GM/R^2$ with $M = \tfrac43\pi R^3\rho$ gives $\rho = 3g/4\pi GR = 5.5\un{g/cm^3}$, twice the density of surface rock. The interior must be much denser than the crust: a core rich in iron, which seismology later confirmed (Problem~\ref{prob:core-shadow}).

\sol{geo}{2} $P = P_0e^{-z/H}$ is halved at $z = H\ln2 = 5.5\un{km}$. Since the pressure at any height is the weight of the air above it, the fraction of the atmosphere above Everest is $e^{-8.85/8} = 33\%$. Climbers on the summit are above two-thirds of the air.

\sol{geo}{3} $v = \sqrt{gh} = \sqrt{9.8\times4000} = 200\un{m/s}$, the speed of a jet airliner. Crossing $17\,000\un{km}$ takes about $24$ hours. (The 1960 Chilean tsunami reached Japan in $22$ hours and killed $140$ people there.) Seismic waves travel at kilometres per second and warnings at the speed of light, so distant coasts can be warned hours in advance.

\sol{geo}{4} $\sigma T_e^4 = S(1-A)/4$. Earth: $T_e = (1361\times0.7/4\sigma)^{1/4} = 255\un{K}$. Venus: $T_e = (2600\times0.23/4\sigma)^{1/4} = 227\un{K}$ --- \emph{colder} than Earth, because its clouds reflect most of the sunlight. Yet its surface is at $737\un{K}$. The difference of $500\un{K}$ is all greenhouse effect, from an atmosphere of $90$ bar of CO$_2$.

\sol{geo}{5} At the top of the atmosphere, only the layer radiates to space, so $\sigma T_a^4 = S(1-A)/4 = \sigma T_e^4$. The surface receives sunlight plus the layer's downward radiation and emits $\sigma T_s^4$, so $\sigma T_s^4 = \sigma T_e^4 + \sigma T_a^4 = 2\sigma T_e^4$, and $T_s = 2^{1/4}T_e = 303\un{K}$. The real Earth is cooler, at $288\un{K}$, because the atmosphere is transparent in an infrared ``window'' near $10\un{\mu m}$ through which the surface radiates directly to space, and because convection and evaporation carry heat from the surface up to where it is radiated more easily.

\sol{geo}{6} $F = \sigma T^4$ gives $\Delta F = 4\sigma T^3\Delta T$, so $\Delta T = 3.7/(4\sigma(255)^3) = 1.0\un{K}$. With $f = 0.5$, $\Delta T = 2\un{K}$. Including ice--albedo and cloud feedbacks, the best current estimate of the warming from doubling CO$_2$ is about $3\un{K}$, likely between $2.5$ and $4\un{K}$.

\sol{geo}{7} The difference between the pulls at the near side and the centre is $GM/(r-R)^2 - GM/r^2\approx2GMR/r^3$. The ratio of lunar to solar tides is $(M_{\text{Moon}}/M_{\text{Sun}})(r_{\text{Sun}}/r_{\text{Moon}})^3 = (3.7\times10^{-8})(389)^3 = 2.2$. The pull itself falls as $1/r^2$, and the Sun's is $180$ times stronger. But the tide falls as $1/r^3$, and the extra factor of $r$ favours the much closer Moon by another factor of $389$.

\sol{geo}{8} For a parcel rising adiabatically, the first law per unit mass gives $c_p\,\dd T = \dd P/\rho$. With $\dd P = -\rho g\,\dd z$, this is $c_p\,\dd T = -g\,\dd z$, so $\Gamma = g/c_p = 9.8\un{K/km}$. The observed average, about $6.5\un{K/km}$, is lower because rising air condenses water vapour, releasing latent heat.

\sol{geo}{9} At $30^\circ$, $f = 2\Omega\sin30^\circ = 7.3\times10^{-5}\un{s^{-1}}$. Bathtub: $\mathrm{Ro} = U/fL = 0.1/(7.3\times10^{-5}\times1)\approx1400$. Hurricane: $30/(7.3\times10^{-5}\times5\times10^5)\approx0.8$. Coriolis forces dominate the hurricane and are negligible in the bathtub, where the direction of swirl is set by leftover motion and the tub's shape. (In careful experiments, with the water left to settle for a day, the Coriolis effect can just be seen.)

\sol{geo}{10} A straight chord between two surface points separated by angle $\Delta$ passes the centre at a distance $R\cos(\Delta/2)$. Setting this equal to the core radius at the shadow edge, $R_c = 6371\cos52^\circ = 3920\un{km}$, within $13\%$ of the true value.

\sol{geo}{11} The Moon's orbital angular momentum is $L_{\text{orb}} = M_m\sqrt{GM_ea} = 2.9\times10^{34}\un{kg\,m^2/s}$, and it grows as $\dd L_{\text{orb}}/\dd t = (L_{\text{orb}}/2a)\,\dd a/\dd t = 4.5\times10^{16}\un{N\,m}$. The Earth loses spin at that rate: $I\,\dd\omega/\dd t = -\dd L_{\text{orb}}/\dd t$, with $I = 0.33M_eR^2 = 8.0\times10^{37}\un{kg\,m^2}$. Since the day is $T = 2\pi/\omega$, $\dd T/\dd t = (T^2/2\pi I)\,\dd L_{\text{orb}}/\dd t = 6.7\times10^{-13}$, or $2.1\un{ms}$ per century --- in excellent agreement with the eclipse records. Over $100$ million years, at the same rate, the day would have lengthened by $35$ minutes: dinosaurs had days of about $23.4$ hours, as growth bands in fossil corals and shells confirm.

\sol{geo}{12} The absorbed sunlight is $(1-A(T))S/4$: $119\un{W/m^2}$ below $250\un{K}$, $238\un{W/m^2}$ above $280\un{K}$, and rising linearly in between. The emitted curve, $0.6\sigma T^4$, is a smooth rising curve. They cross three times: at $243\un{K}$ on the icy branch, at $258\un{K}$ on the slope, and at $289\un{K}$ on the warm branch. At the outer two, a small warming makes emission exceed absorption, and the planet cools back: stable. At the middle one, a small warming melts ice, lowers the albedo, and runs away to the warm state; a small cooling runs away to the snowball. A dip in sunlight or greenhouse gas can push the climate past the middle point, into the snowball. To escape, the emission curve must drop until the snowball equilibrium disappears: absorption must exceed emission even at $250\un{K}$, which needs $\epsilon<119/(\sigma\cdot250^4) = 0.54$. Volcanoes keep emitting CO$_2$, and on a frozen planet there is no rain to remove it by weathering rocks, so over millions of years it builds up. When the threshold is passed, the Earth jumps abruptly to the warm branch --- now with a strong greenhouse effect, at around $300\un{K}$ --- and the ice melts in a geological instant. Geologists find the evidence: glacial deposits in the tropics, capped by thick layers of carbonate rock laid down in the hot aftermath. The path out is not the path in: hysteresis.
